\chapter[\emj{🔁}~Detección de Ciclos en Grafos]{\emj{🔁}~Detección de Ciclos en Grafos}

\begin{dsaquees}

Caminar por una ciudad 🚶 y de repente darte cuenta de que ya pasaste
por esa esquina: diste una vuelta completa. Eso es un ciclo.

Detectarlo depende del tipo de calle:
\begin{itemize}
\item Calles de doble sentido (grafo no dirigido): hay ciclo si llegas
      a un lugar ya visitado que NO es de donde acabas de venir.
\item Calles de un solo sentido (grafo dirigido): hay ciclo si vuelves
      a pisar un lugar que todavía está "en tu camino actual", no uno
      que ya terminaste de explorar.
\end{itemize}

\end{dsaquees}

\begin{dsaotraforma}

Detectar ciclos es fundamental: un ciclo cambia por completo qué
algoritmos aplican (por ejemplo, el ordenamiento topológico solo existe
si NO hay ciclos).

En grafos \textbf{no dirigidos} usamos DFS guardando el "padre": si
encontramos un vecino ya visitado distinto del padre, hay ciclo.
Alternativa: Union-Find (si dos extremos de una arista ya están unidos,
esa arista cierra un ciclo).

En grafos \textbf{dirigidos} usamos DFS con 3 estados (colores):
Blanco (sin visitar), Gris (en la pila de recursión actual) y Negro
(terminado). Si durante el DFS tocamos un nodo Gris, hay ciclo. El
algoritmo de Kahn también sirve: si el topo-sort no procesa todos los
nodos, hay ciclo.

\end{dsaotraforma}

\begin{dsadiagrama}
\begin{verbatim}
NO DIRIGIDO (padre):        DIRIGIDO (3 colores):

  A ─── B                     A ──> B ──> C
  │     │                     ^           │
  D ─── C                     └───────────┘

DFS: A->B->C->D->A            DFS desde A:
Al ver A de nuevo (no es      A(gris) B(gris) C(gris)
padre de D) => CICLO          C apunta a A que esta GRIS
                              => CICLO

Estados: Blanco=0  Gris=1 (en recursion)  Negro=2 (listo)
\end{verbatim}
\end{dsadiagrama}

\begin{dsacomofunciona}
\begin{verbatim}
TRAZA 3 COLORES (dirigido)   grafo: 0→1→2→0

Paso │ Nodo │ Acción              │ Estados (0,1,2)
─────┼──────┼─────────────────────┼──────────────────
  1  │  0   │ entra → GRIS        │ (G, B, B)
  2  │  1   │ entra → GRIS        │ (G, G, B)
  3  │  2   │ entra → GRIS        │ (G, G, G)
  4  │  2   │ vecino 0 es GRIS →  │  ¡CICLO detectado!
                arista de retorno

Blanco(B)=sin visitar  Gris(G)=en la pila actual  Negro(N)=terminado
Ver un GRIS = arista hacia un ancestro = ciclo.
Ver un NEGRO ≠ ciclo (ya cerró, es otra rama).

NO DIRIGIDO (con padre):   0—1—2   y   2—0
  DFS 0(pad=-1) → 1(pad=0) → 2(pad=1)
  desde 2 veo vecino 0: visitado y 0≠padre(1) → CICLO
\end{verbatim}
\end{dsacomofunciona}

\begin{lstlisting}
NO DIRIGIDO (DFS con padre)
- Recorre con DFS. Al visitar un vecino v de u:
  * Si v no visitado -> DFS(v, padre=u).
  * Si v visitado y v != padre -> hay ciclo.

DIRIGIDO (DFS con 3 colores)
- color[u] = GRIS al entrar, NEGRO al salir.
- Si un vecino v es GRIS -> hay ciclo (arista hacia atrás).

PSEUDOCÓDIGO (dirigido, 3 colores)
──────────────────────────────────────────────────────────────

función tieneCicloDirigido(u):
    color[u] = GRIS
    PARA cada vecino v de u:
        SI color[v] == GRIS: return true          // ciclo
        SI color[v] == BLANCO Y tieneCicloDirigido(v): return true
    color[u] = NEGRO
    return false

CÓDIGO C++ (no dirigido, con padre)
──────────────────────────────────────────────────────────────

#include <vector>
using namespace std;

bool dfsUndirected(int u, int parent, vector<vector<int>>& adj,
                   vector<bool>& vis) {
    vis[u] = true;
    for (int v : adj[u]) {
        if (!vis[v]) {
            if (dfsUndirected(v, u, adj, vis)) return true;
        } else if (v != parent) {
            return true;              // visitado y no es el padre
        }
    }
    return false;
}

CÓDIGO C++ (dirigido, 3 colores)
──────────────────────────────────────────────────────────────

bool dfsDirected(int u, vector<vector<int>>& adj, vector<int>& color) {
    color[u] = 1;                     // GRIS (en recursión)
    for (int v : adj[u]) {
        if (color[v] == 1) return true;              // ciclo
        if (color[v] == 0 && dfsDirected(v, adj, color)) return true;
    }
    color[u] = 2;                     // NEGRO (terminado)
    return false;
}

COMPLEJIDAD (Big-O)
──────────────────────────────────────────────────────────────

  Método            │ Tiempo      │ Espacio
  ──────────────────┼─────────────┼──────────────
  DFS (ambos)       │ O(V + E)    │ O(V)
  Union-Find (undir)│ O(E α(V))   │ O(V)

* α = función inversa de Ackermann, prácticamente constante.
\end{lstlisting}

\begin{dsaentrevistas}

✅ Elige el método según el grafo: no dirigido -> padre o Union-Find;
   dirigido -> 3 colores o Kahn. Confundirlos es un error típico.

💡 "Course Schedule" (LeetCode 207): "¿es posible terminar el plan?" =
   "¿el grafo de prerrequisitos NO tiene ciclos?". Ciclo dirigido.

⚠️ El truco del "padre" solo aplica a NO dirigidos. En dirigidos un
   nodo NEGRO (ya terminado) no implica ciclo; solo un GRIS lo implica.

🔑 Detección de deadlocks: en sistemas operativos, un ciclo en el grafo
   de "espera de recursos" = interbloqueo. Buen dato para mencionar.

\end{dsaentrevistas}

\begin{dsaresumen}

• No dirigido: DFS con padre, o Union-Find (arista que une lo ya unido).
• Dirigido: DFS con 3 colores; un vecino GRIS = ciclo.
• Alternativa dirigida: Kahn no procesa todos los nodos = ciclo.
• Todo en O(V + E).
• Vital para topo-sort, detección de deadlocks y validar DAGs.
════════════════════════════════════════════════════════════════

\end{dsaresumen}
